\documentclass[11pt]{article}
\usepackage[a4paper,margin=1.9cm,top=2.4cm,bottom=2.2cm]{geometry}
\usepackage{amsmath,amssymb,amsfonts,fontspec,microtype,fancyhdr,lastpage,enumitem,needspace}
\usepackage[most]{tcolorbox}
\usepackage{xcolor}
\definecolor{ink}{HTML}{1F3A5F}\definecolor{soft}{HTML}{EEF3F9}\definecolor{ans}{HTML}{F3F8F1}\definecolor{ansb}{HTML}{3C7A3C}
\pagestyle{fancy}\fancyhf{}
\lhead{\small\color{ink}\textbf{Linear Algebra I} \textperiodcentered\ Week 4 --- Row Reduction, Solving Linear Systems \& Bases from RREF}
\rhead{}\cfoot{\small Mr Malek Thiab \textperiodcentered\ Worksheet with Worked Answers \textperiodcentered\ Page \thepage\ of \pageref{LastPage}}
\renewcommand{\headrulewidth}{0.4pt}
\setlength{\parindent}{0pt}\setlength{\parskip}{4pt}
\newcommand{\lv}[1]{\textsc{\small #1}}
\begin{document}
{\Large\bfseries\color{ink} Linear Algebra I --- Week 4}\\[2pt]
{\large\bfseries\color{ink} Row Reduction, Solving Linear Systems \& Bases from RREF}\\[3pt]
{\small Name: \underline{\hspace{5cm}}\hfill Date: \underline{\hspace{3cm}}\qquad All work \textbf{by hand} --- no calculator, no software.}
\vspace{4pt}\hrule\vspace{8pt}

\begin{tcolorbox}[colback=soft,colframe=ink,title={\bfseries Key Facts}]
\textbf{Elementary row operations.} (i) swap two rows $R_i\leftrightarrow R_j$; (ii) scale a row $R_i\to cR_i$ with $c\neq0$; (iii) add a multiple of one row to another $R_i\to R_i+\lambda R_j$. They do not change the solution set of a system, and they can be undone, so they do not change the span of the rows.

\textbf{REF and RREF.} A matrix is in \emph{row echelon form} (REF) if the zero rows are at the bottom and the first nonzero entry (the \emph{leading entry}) of each row lies strictly to the right of the leading entry of the row above. It is in \emph{reduced} row echelon form (RREF) if moreover every leading entry is $1$ and is the only nonzero entry in its column. The RREF of a matrix is unique. A \emph{pivot column} is a column containing a leading entry.

\textbf{Solving $A\mathbf x=\mathbf b$.} Reduce $\bigl(A\mid\mathbf b\bigr)$ to RREF. If a row reads $(0\ \cdots\ 0\mid1)$ the system is \emph{inconsistent}. Otherwise variables in pivot columns are \emph{basic}, the others are \emph{free}. No free variable: a unique solution. At least one free variable: infinitely many solutions, written as $\mathbf p+s\,\mathbf v_1+t\,\mathbf v_2+\cdots$ (parametric vector form).

\textbf{Homogeneous systems.} $A\mathbf x=\mathbf 0$ is always consistent ($\mathbf x=\mathbf 0$). Its solution set is a subspace, the span of the vectors multiplying the free parameters. If there are more unknowns than equations there is a free variable, hence a nonzero solution. If $\mathbf p$ solves $A\mathbf x=\mathbf b$, every solution is $\mathbf p+\mathbf h$ with $A\mathbf h=\mathbf 0$.

\textbf{Columns, rows and RREF.} The columns of $A$ at the pivot positions of its RREF form a basis of $\operatorname{col}(A)$ (take them from the original $A$, not from the RREF). The nonzero rows of the RREF form a basis of $\operatorname{row}(A)$. The entries of a non-pivot column of the RREF are the coefficients expressing that column of $A$ through the pivot columns. Vectors $v_1,\dots,v_k$ are independent exactly when the matrix with these columns has a pivot in every column.
\end{tcolorbox}

\newpage
\section*{\large\color{ink} Worked Examples}
\begin{tcolorbox}[breakable,colback=white,colframe=ink!60,title={\bfseries Example 1 \textperiodcentered\ Reducing a matrix to RREF}]Reduce $\begin{pmatrix}1 & 2 & 1 \\ 2 & 5 & 3 \\ 1 & 3 & 4\end{pmatrix}$ to RREF.
\par\smallskip
\textbf{Solution.} \begin{enumerate}[leftmargin=1.6em,itemsep=1pt,topsep=2pt,label=\arabic*.]\item Go column by column. Column 1: the top entry is already $1$, so use it as the pivot and clear the entries below.
\item Apply $R_2\to R_2-2R_1,\ R_3\to R_3-R_1$:\par\smallskip $\displaystyle \begin{pmatrix}1 & 2 & 1 \\ 0 & 1 & 1 \\ 0 & 1 & 3\end{pmatrix}$
\item Apply $R_1\to R_1-2R_2,\ R_3\to R_3-R_2$:\par\smallskip $\displaystyle \begin{pmatrix}1 & 0 & -1 \\ 0 & 1 & 1 \\ 0 & 0 & 2\end{pmatrix}$
\item Apply $R_3\to \dfrac{1}{2} R_3$:\par\smallskip $\displaystyle \begin{pmatrix}1 & 0 & -1 \\ 0 & 1 & 1 \\ 0 & 0 & 1\end{pmatrix}$
\item Apply $R_1\to R_1+R_3,\ R_2\to R_2-R_3$:\par\smallskip $\displaystyle \begin{pmatrix}1 & 0 & 0 \\ 0 & 1 & 0 \\ 0 & 0 & 1\end{pmatrix}$
\item Every column has a pivot; the result is the identity matrix, which is in RREF.
\end{enumerate}\textbf{Answer:} RREF $=\begin{pmatrix}1 & 0 & 0 \\ 0 & 1 & 0 \\ 0 & 0 & 1\end{pmatrix}$.
\end{tcolorbox}
\begin{tcolorbox}[breakable,colback=white,colframe=ink!60,title={\bfseries Example 2 \textperiodcentered\ A system with a unique solution}]Solve $\begin{cases}x+y=3\\ 2x-y=0\end{cases}$
\par\smallskip
\textbf{Solution.} \begin{enumerate}[leftmargin=1.6em,itemsep=1pt,topsep=2pt,label=\arabic*.]\item Write the system as an augmented matrix, one row per equation, with a bar before the right-hand side:\par\smallskip $\displaystyle \left(\begin{array}{cc|c}1 & 1 & 3 \\ 2 & -1 & 0\end{array}\right)$
\item Apply $R_2\to R_2-2R_1$:\par\smallskip $\displaystyle \left(\begin{array}{cc|c}1 & 1 & 3 \\ 0 & -3 & -6\end{array}\right)$
\item Apply $R_2\to \dfrac{1}{-3} R_2$:\par\smallskip $\displaystyle \left(\begin{array}{cc|c}1 & 1 & 3 \\ 0 & 1 & 2\end{array}\right)$
\item Apply $R_1\to R_1-R_2$:\par\smallskip $\displaystyle \left(\begin{array}{cc|c}1 & 0 & 1 \\ 0 & 1 & 2\end{array}\right)$
\item Every column of the coefficient part holds a pivot, so there are no free variables and the solution is unique.
\item Read each nonzero row as an equation for its basic variable:\par\smallskip $\displaystyle\begin{aligned}x&=1\\ y&=2\end{aligned}$
\item Check by substituting $(x,y)=(1,2)$ into the original equations: the left sides give $(3,0)$, the right-hand sides \checkmark.
\end{enumerate}\textbf{Answer:} Unique solution: $(x,y)=(1,2)$.
\end{tcolorbox}
\begin{tcolorbox}[breakable,colback=white,colframe=ink!60,title={\bfseries Example 3 \textperiodcentered\ A system with infinitely many solutions}]Solve $\begin{cases}x+y+z=3\\ 2x+y=3\end{cases}$
\par\smallskip
\textbf{Solution.} \begin{enumerate}[leftmargin=1.6em,itemsep=1pt,topsep=2pt,label=\arabic*.]\item Write the system as an augmented matrix, one row per equation, with a bar before the right-hand side:\par\smallskip $\displaystyle \left(\begin{array}{ccc|c}1 & 1 & 1 & 3 \\ 2 & 1 & 0 & 3\end{array}\right)$
\item Apply $R_2\to R_2-2R_1$:\par\smallskip $\displaystyle \left(\begin{array}{ccc|c}1 & 1 & 1 & 3 \\ 0 & -1 & -2 & -3\end{array}\right)$
\item Apply $R_2\to \dfrac{1}{-1} R_2$:\par\smallskip $\displaystyle \left(\begin{array}{ccc|c}1 & 1 & 1 & 3 \\ 0 & 1 & 2 & 3\end{array}\right)$
\item Apply $R_1\to R_1-R_2$:\par\smallskip $\displaystyle \left(\begin{array}{ccc|c}1 & 0 & -1 & 0 \\ 0 & 1 & 2 & 3\end{array}\right)$
\item Pivot columns are the columns of the basic variables: $x$, $y$. The remaining variable(s) $z$ are \emph{free}: call them $z=s$.
\item Read each nonzero row as an equation for its basic variable:\par\smallskip $\displaystyle\begin{aligned}x&=s\\ y&=3-2s\\ z&=s\end{aligned}$
\item Write this as a vector: $\displaystyle \mathbf x=\begin{pmatrix}0\\3\\0\end{pmatrix}+s\begin{pmatrix}1\\-2\\1\end{pmatrix}$ (the parameters $s$ may be any numbers).
\item Check with $s=1$: $\mathbf x=\begin{pmatrix}1\\1\\1\end{pmatrix}$, and substituting into the original equations gives the right-hand sides $(3,3)$ \checkmark.
\end{enumerate}\textbf{Answer:} Infinitely many solutions: $\mathbf x=\begin{pmatrix}0\\3\\0\end{pmatrix}+s\begin{pmatrix}1\\-2\\1\end{pmatrix}$.
\end{tcolorbox}
\begin{tcolorbox}[breakable,colback=white,colframe=ink!60,title={\bfseries Example 4 \textperiodcentered\ An inconsistent system}]Solve $\begin{cases}x+y+2z=1\\ 2x+2y+5z=3\\ x+y=2\end{cases}$
\par\smallskip
\textbf{Solution.} \begin{enumerate}[leftmargin=1.6em,itemsep=1pt,topsep=2pt,label=\arabic*.]\item Write the system as an augmented matrix, one row per equation, with a bar before the right-hand side:\par\smallskip $\displaystyle \left(\begin{array}{ccc|c}1 & 1 & 2 & 1 \\ 2 & 2 & 5 & 3 \\ 1 & 1 & 0 & 2\end{array}\right)$
\item Apply $R_2\to R_2-2R_1,\ R_3\to R_3-R_1$:\par\smallskip $\displaystyle \left(\begin{array}{ccc|c}1 & 1 & 2 & 1 \\ 0 & 0 & 1 & 1 \\ 0 & 0 & -2 & 1\end{array}\right)$
\item Apply $R_1\to R_1-2R_2,\ R_3\to R_3+2R_2$:\par\smallskip $\displaystyle \left(\begin{array}{ccc|c}1 & 1 & 0 & -1 \\ 0 & 0 & 1 & 1 \\ 0 & 0 & 0 & 3\end{array}\right)$
\item Row 3 reads $0=1$, a contradiction. No numbers $x,y,z$ can satisfy it, so the system is \emph{inconsistent}.
\end{enumerate}\textbf{Answer:} No solution (inconsistent).
\end{tcolorbox}
\begin{tcolorbox}[breakable,colback=white,colframe=ink!60,title={\bfseries Example 5 \textperiodcentered\ A homogeneous system}]Solve $\begin{cases}x+y+2w=0\\ 2x+2y+z+5w=0\end{cases}$ and write the solution set as a span.
\par\smallskip
\textbf{Solution.} \begin{enumerate}[leftmargin=1.6em,itemsep=1pt,topsep=2pt,label=\arabic*.]\item Write the system as an augmented matrix, one row per equation, with a bar before the right-hand side:\par\smallskip $\displaystyle \left(\begin{array}{cccc|c}1 & 1 & 0 & 2 & 0 \\ 2 & 2 & 1 & 5 & 0\end{array}\right)$
\item Apply $R_2\to R_2-2R_1$:\par\smallskip $\displaystyle \left(\begin{array}{cccc|c}1 & 1 & 0 & 2 & 0 \\ 0 & 0 & 1 & 1 & 0\end{array}\right)$
\item Pivot columns are the columns of the basic variables: $x$, $z$. The remaining variable(s) $y$, $w$ are \emph{free}: call them $y=s$, $w=t$.
\item Read each nonzero row as an equation for its basic variable:\par\smallskip $\displaystyle\begin{aligned}x&=-s-2t\\ y&=s\\ z&=-t\\ w&=t\end{aligned}$
\item Write this as a vector: $\displaystyle \mathbf x=s\begin{pmatrix}-1\\1\\0\\0\end{pmatrix}+t\begin{pmatrix}-2\\0\\-1\\1\end{pmatrix}$ (the parameters $s$, $t$ may be any numbers).
\item Check with $s=1$, $t=1$: $\mathbf x=\begin{pmatrix}-3\\1\\-1\\1\end{pmatrix}$, and substituting into the original equations gives the right-hand sides $(0,0)$ \checkmark.
\end{enumerate}\textbf{Answer:} Infinitely many solutions: $\mathbf x=s\begin{pmatrix}-1\\1\\0\\0\end{pmatrix}+t\begin{pmatrix}-2\\0\\-1\\1\end{pmatrix}$. The solution set is $\operatorname{span}$ of the two vectors above.
\end{tcolorbox}
\begin{tcolorbox}[breakable,colback=white,colframe=ink!60,title={\bfseries Example 6 \textperiodcentered\ Column space basis and column relations}]For $A=\begin{pmatrix}1 & 3 & 0 & 2 \\ 2 & 6 & 1 & 5 \\ 1 & 3 & 1 & 3\end{pmatrix}$ find a basis of $\operatorname{col}(A)$ and express the other columns using it.
\par\smallskip
\textbf{Solution.} \begin{enumerate}[leftmargin=1.6em,itemsep=1pt,topsep=2pt,label=\arabic*.]\item Reduce $A$ to RREF.
\item Apply $R_2\to R_2-2R_1,\ R_3\to R_3-R_1$:\par\smallskip $\displaystyle \begin{pmatrix}1 & 3 & 0 & 2 \\ 0 & 0 & 1 & 1 \\ 0 & 0 & 1 & 1\end{pmatrix}$
\item Apply $R_3\to R_3-R_2$:\par\smallskip $\displaystyle \begin{pmatrix}1 & 3 & 0 & 2 \\ 0 & 0 & 1 & 1 \\ 0 & 0 & 0 & 0\end{pmatrix}$
\item Pivot columns: 1,3. By the pivot-column rule, a basis of $\operatorname{col}(A)$ is formed by columns 1,3 \emph{of the original $A$}: $\left\{\begin{pmatrix}1\\2\\1\end{pmatrix},\ \begin{pmatrix}0\\1\\1\end{pmatrix}\right\}$.
\item The non-pivot columns of the RREF give the coefficients: column 2 of $R$ is $\begin{pmatrix}3\\0\end{pmatrix}$, so $\text{col}_2=3\,\text{col}_1$; column 4 of $R$ is $\begin{pmatrix}2\\1\end{pmatrix}$, so $\text{col}_4=2\,\text{col}_1+1\,\text{col}_3$.
\end{enumerate}\textbf{Answer:} Basis $\{\begin{pmatrix}1\\2\\1\end{pmatrix},\begin{pmatrix}0\\1\\1\end{pmatrix}\}$.
\end{tcolorbox}
\begin{tcolorbox}[breakable,colback=white,colframe=ink!60,title={\bfseries Example 7 \textperiodcentered\ A system over $\mathbb{C}$}]Solve $\begin{cases}x+iy=1+i\\ ix+2y=i\end{cases}$ over $\mathbb{C}$.
\par\smallskip
\textbf{Solution.} \begin{enumerate}[leftmargin=1.6em,itemsep=1pt,topsep=2pt,label=\arabic*.]\item Write the system as an augmented matrix, one row per equation, with a bar before the right-hand side:\par\smallskip $\displaystyle \left(\begin{array}{cc|c}1 & i & 1+i \\ i & 2 & i\end{array}\right)$
\item Apply $R_2\to R_2-iR_1$:\par\smallskip $\displaystyle \left(\begin{array}{cc|c}1 & i & 1+i \\ 0 & 3 & 1\end{array}\right)$
\item Apply $R_2\to \dfrac{1}{3} R_2$:\par\smallskip $\displaystyle \left(\begin{array}{cc|c}1 & i & 1+i \\ 0 & 1 & \dfrac{1}{3}\end{array}\right)$
\item Apply $R_1\to R_1-iR_2$:\par\smallskip $\displaystyle \left(\begin{array}{cc|c}1 & 0 & 1+\dfrac{2}{3}i \\ 0 & 1 & \dfrac{1}{3}\end{array}\right)$
\item Every column of the coefficient part holds a pivot, so there are no free variables and the solution is unique.
\item Read each nonzero row as an equation for its basic variable:\par\smallskip $\displaystyle\begin{aligned}x&=1+\dfrac{2}{3}i\\ y&=\dfrac{1}{3}\end{aligned}$
\item Check by substituting $(x,y)=(1+\dfrac{2}{3}i,\dfrac{1}{3})$ into the original equations: the left sides give $(1+i,i)$, the right-hand sides \checkmark.
\end{enumerate}\textbf{Answer:} Unique solution: $(x,y)=(1+\dfrac{2}{3}i,\dfrac{1}{3})$.
\end{tcolorbox}
\begin{tcolorbox}[breakable,colback=white,colframe=ink!60,title={\bfseries Example 8 \textperiodcentered\ Proof: row reduction does not change the solutions}]Prove: swapping two equations, scaling an equation by $c\neq0$, or adding a multiple of one equation to another does not change the solution set.
\par\smallskip
\textbf{Solution.} \begin{enumerate}[leftmargin=1.6em,itemsep=1pt,topsep=2pt,label=\arabic*.]\item Swapping: the system is the same list of equations in a different order, and a vector solves a list of equations regardless of order.
\item Scaling $E_i$ by $c\neq0$: if $E_i(\mathbf x)=\beta_i$ then $(cE_i)(\mathbf x)=c\beta_i$. Conversely, if $(cE_i)(\mathbf x)=c\beta_i$, multiply both sides by $c^{-1}$ (possible because $c\neq0$ in a field) to recover $E_i(\mathbf x)=\beta_i$.
\item Adding $\lambda E_j$ to $E_i$: if $\mathbf x$ solves both $E_i,E_j$ it solves $E_i+\lambda E_j$. Conversely, if $\mathbf x$ solves $E_j$ and $E_i+\lambda E_j$, then $E_i(\mathbf x)=(E_i+\lambda E_j)(\mathbf x)-\lambda E_j(\mathbf x)=\beta_i$.
\item Each operation is undone by an operation of the same kind, so solutions are preserved in both directions.
\end{enumerate}\textbf{Answer:} The three operations preserve the solution set, so a system and its RREF have the same solutions.
\end{tcolorbox}
\newpage
\section*{\large\color{ink} Practice Problems}
\Needspace{16\baselineskip}\subsection*{\normalsize\color{ink} Part A --- Echelon Forms and Row Reduction}
\Needspace{9\baselineskip}\begin{minipage}{\linewidth}\textbf{1.} \hfill{\scriptsize\color{ink}[Foundational]}\par\nopagebreak Classify each matrix as: in RREF, in REF but not RREF, or neither. For every matrix that is not in RREF, name the condition it violates. $$\text{(a) }\begin{pmatrix}1 & 0 & 2 \\ 0 & 1 & -1 \\ 0 & 0 & 0\end{pmatrix}\quad \text{(b) }\begin{pmatrix}1 & 0 & 2 \\ 0 & 1 & 0 \\ 0 & 0 & 1\end{pmatrix}\quad\text{(c) }\begin{pmatrix}0 & 1 & 3 \\ 1 & 0 & 2 \\ 0 & 0 & 0\end{pmatrix}\quad\text{(d) }\begin{pmatrix}1 & 3 & 0 & 2 \\ 0 & 0 & 1 & -1 \\ 0 & 0 & 0 & 0\end{pmatrix}$$
\end{minipage}\par\nopagebreak
\begin{tcolorbox}[breakable,colback=ans,colframe=ansb,boxrule=0.5pt,left=4pt,right=4pt,top=2pt,bottom=2pt]\begin{enumerate}[leftmargin=2.2em,itemsep=1pt,topsep=1pt,label=\textbf{Step \arabic*.}]\item Recall the two lists of conditions. \emph{REF:} zero rows are at the bottom, and the first nonzero entry (the \emph{leading entry}) of each row lies strictly to the right of the leading entry of the row above. \emph{RREF:} REF, plus every leading entry equals $1$, plus each leading entry is the only nonzero entry in its column.
\item (a) Leading entries: row 1 in column 1, row 2 in column 2, row 3 is a zero row at the bottom. They move right, so it is REF. Both leading entries are $1$. Column 1 has only the entry $1$ (others $0$). Column 2 has only the entry $1$. So (a) is in RREF.
\item (b) Leading entries sit in columns $1,2,3$, moving right, so it is REF. All leading entries are $1$. But column 3 contains the leading entry $1$ of row 3 and also the entry $2$ above it in row 1. A leading entry must be the \emph{only} nonzero entry in its column, so (b) is REF but \emph{not} RREF.
\item (c) Row 1 has its leading entry in column 2; row 2 has its leading entry in column 1, which is to the \emph{left} of the one above. The leading entries do not move to the right, so (c) is not even REF.
\item (d) Row 1 leads in column 1, row 2 leads in column 3, row 3 is zero and at the bottom: REF. Leading entries are $1$. Column 1 has no other nonzero entry; column 3 has only the $1$ (the entry above it is $0$; the $3$ and $2$ lie in non-pivot columns, which is allowed). So (d) is in RREF.
\end{enumerate}\textbf{Answer:} (a) RREF. (b) REF only: the entry $2$ above the pivot in column 3 is nonzero. (c) Neither: leading entries do not move to the right. (d) RREF.\end{tcolorbox}
\par\vspace{2pt}
\Needspace{9\baselineskip}\begin{minipage}{\linewidth}\textbf{2.} \hfill{\scriptsize\color{ink}[Foundational]}\par\nopagebreak Use elementary row operations to reduce $\begin{pmatrix}1 & 2 & 1 \\ 2 & 5 & 3 \\ 1 & 3 & 4\end{pmatrix}$ to RREF. Write each operation in the form $R_i\to R_i-cR_j$ and state the pivot columns.
\end{minipage}\par\nopagebreak
\begin{tcolorbox}[breakable,colback=ans,colframe=ansb,boxrule=0.5pt,left=4pt,right=4pt,top=2pt,bottom=2pt]\begin{enumerate}[leftmargin=2.2em,itemsep=1pt,topsep=1pt,label=\textbf{Step \arabic*.}]\item Strategy: take each column in turn from the left, make its pivot equal to $1$, then clear every other entry in that column using the pivot row.
\item Apply $R_2\to R_2-2R_1,\ R_3\to R_3-R_1$:\par\smallskip $\displaystyle \begin{pmatrix}1 & 2 & 1 \\ 0 & 1 & 1 \\ 0 & 1 & 3\end{pmatrix}$
\item Apply $R_1\to R_1-2R_2,\ R_3\to R_3-R_2$:\par\smallskip $\displaystyle \begin{pmatrix}1 & 0 & -1 \\ 0 & 1 & 1 \\ 0 & 0 & 2\end{pmatrix}$
\item Apply $R_3\to \dfrac{1}{2} R_3$:\par\smallskip $\displaystyle \begin{pmatrix}1 & 0 & -1 \\ 0 & 1 & 1 \\ 0 & 0 & 1\end{pmatrix}$
\item Apply $R_1\to R_1+R_3,\ R_2\to R_2-R_3$:\par\smallskip $\displaystyle \begin{pmatrix}1 & 0 & 0 \\ 0 & 1 & 0 \\ 0 & 0 & 1\end{pmatrix}$
\item The matrix is now in RREF. It has a pivot in each of columns $1,2,3$.
\end{enumerate}\textbf{Answer:} RREF $=\begin{pmatrix}1 & 0 & 0 \\ 0 & 1 & 0 \\ 0 & 0 & 1\end{pmatrix}$; pivot columns $1,2,3$.\end{tcolorbox}
\par\vspace{2pt}
\Needspace{9\baselineskip}\begin{minipage}{\linewidth}\textbf{3.} \hfill{\scriptsize\color{ink}[Intermediate]}\par\nopagebreak Reduce $\begin{pmatrix}1 & 2 & -1 & 3 \\ 2 & 4 & 0 & 2 \\ 1 & 2 & 1 & -1\end{pmatrix}$ to RREF. State the pivot columns and explain why column 2 is not a pivot column.
\end{minipage}\par\nopagebreak
\begin{tcolorbox}[breakable,colback=ans,colframe=ansb,boxrule=0.5pt,left=4pt,right=4pt,top=2pt,bottom=2pt]\begin{enumerate}[leftmargin=2.2em,itemsep=1pt,topsep=1pt,label=\textbf{Step \arabic*.}]\item Work column by column. Column 1 already has a $1$ in the top position, so use it to clear below.
\item Apply $R_2\to R_2-2R_1,\ R_3\to R_3-R_1$:\par\smallskip $\displaystyle \begin{pmatrix}1 & 2 & -1 & 3 \\ 0 & 0 & 2 & -4 \\ 0 & 0 & 2 & -4\end{pmatrix}$
\item Apply $R_2\to \dfrac{1}{2} R_2$:\par\smallskip $\displaystyle \begin{pmatrix}1 & 2 & -1 & 3 \\ 0 & 0 & 1 & -2 \\ 0 & 0 & 2 & -4\end{pmatrix}$
\item Apply $R_1\to R_1+R_2,\ R_3\to R_3-2R_2$:\par\smallskip $\displaystyle \begin{pmatrix}1 & 2 & 0 & 1 \\ 0 & 0 & 1 & -2 \\ 0 & 0 & 0 & 0\end{pmatrix}$
\item After clearing column 1, column 2 below row 1 is entirely $0$, so there is no pivot in column 2 (nothing available to turn into a leading entry). Move on to column 3: it has a nonzero entry in row 2, which becomes the pivot.
\item The matrix is in RREF with pivot columns $1$ and $3$.
\end{enumerate}\textbf{Answer:} RREF $=\begin{pmatrix}1 & 2 & 0 & 1 \\ 0 & 0 & 1 & -2 \\ 0 & 0 & 0 & 0\end{pmatrix}$. Pivot columns $1,3$; column 2 has no pivot because, once column 1 is cleared, every entry of column 2 in rows $2$ and $3$ is $0$.\end{tcolorbox}
\par\vspace{2pt}
\Needspace{9\baselineskip}\begin{minipage}{\linewidth}\textbf{4.} \hfill{\scriptsize\color{ink}[Intermediate]}\par\nopagebreak Reduce $\begin{pmatrix}0 & 2 & 4 \\ 1 & 1 & 1 \\ 2 & 0 & -2\end{pmatrix}$ to RREF. A row swap is needed: say why, and state the pivot columns.
\end{minipage}\par\nopagebreak
\begin{tcolorbox}[breakable,colback=ans,colframe=ansb,boxrule=0.5pt,left=4pt,right=4pt,top=2pt,bottom=2pt]\begin{enumerate}[leftmargin=2.2em,itemsep=1pt,topsep=1pt,label=\textbf{Step \arabic*.}]\item The entry in position $(1,1)$ is $0$, so it cannot serve as a pivot. We must bring up a row whose first entry is nonzero: swap row 1 with row 2 (or row 3). The swap is an allowed elementary row operation and does not change the solution set of a system.
\item Apply $R_1\leftrightarrow R_2$:\par\smallskip $\displaystyle \begin{pmatrix}1 & 1 & 1 \\ 0 & 2 & 4 \\ 2 & 0 & -2\end{pmatrix}$
\item Apply $R_3\to R_3-2R_1$:\par\smallskip $\displaystyle \begin{pmatrix}1 & 1 & 1 \\ 0 & 2 & 4 \\ 0 & -2 & -4\end{pmatrix}$
\item Apply $R_2\to \dfrac{1}{2} R_2$:\par\smallskip $\displaystyle \begin{pmatrix}1 & 1 & 1 \\ 0 & 1 & 2 \\ 0 & -2 & -4\end{pmatrix}$
\item Apply $R_1\to R_1-R_2,\ R_3\to R_3+2R_2$:\par\smallskip $\displaystyle \begin{pmatrix}1 & 0 & -1 \\ 0 & 1 & 2 \\ 0 & 0 & 0\end{pmatrix}$
\item The matrix is now in RREF.
\end{enumerate}\textbf{Answer:} RREF $=\begin{pmatrix}1 & 0 & -1 \\ 0 & 1 & 2 \\ 0 & 0 & 0\end{pmatrix}$; pivot columns 1,2.\end{tcolorbox}
\par\vspace{2pt}
\Needspace{9\baselineskip}\begin{minipage}{\linewidth}\textbf{5.} \hfill{\scriptsize\color{ink}[Challenge]}\par\nopagebreak Over $\mathbb{C}$, reduce $\begin{pmatrix}1 & i & 0 \\ 2 & 2i & 1 \\ 1 & i & 1\end{pmatrix}$ to RREF. State the pivot columns.
\end{minipage}\par\nopagebreak
\begin{tcolorbox}[breakable,colback=ans,colframe=ansb,boxrule=0.5pt,left=4pt,right=4pt,top=2pt,bottom=2pt]\begin{enumerate}[leftmargin=2.2em,itemsep=1pt,topsep=1pt,label=\textbf{Step \arabic*.}]\item Arithmetic over $\mathbb{C}$ follows the same rules; the only new fact needed is $i^{2}=-1$. Column 1 has a $1$ at the top, so clear below it. The multiples of $i$ in column 2 will be cleared automatically along with column 1.
\item Apply $R_2\to R_2-2R_1,\ R_3\to R_3-R_1$:\par\smallskip $\displaystyle \begin{pmatrix}1 & i & 0 \\ 0 & 0 & 1 \\ 0 & 0 & 1\end{pmatrix}$
\item Apply $R_3\to R_3-R_2$:\par\smallskip $\displaystyle \begin{pmatrix}1 & i & 0 \\ 0 & 0 & 1 \\ 0 & 0 & 0\end{pmatrix}$
\item Column 2 is zero below row 1 (it is $i$ times column 1), so it holds no pivot. Column 3 gives the second pivot.
\end{enumerate}\textbf{Answer:} RREF $=\begin{pmatrix}1 & i & 0 \\ 0 & 0 & 1 \\ 0 & 0 & 0\end{pmatrix}$; pivot columns $1,3$.\end{tcolorbox}
\par\vspace{2pt}
\Needspace{9\baselineskip}\begin{minipage}{\linewidth}\textbf{6.} \hfill{\scriptsize\color{ink}[Challenge]}\par\nopagebreak (Why row reduction works.) Let $S'$ be obtained from a system $S$ by one elementary row operation $R_i\to R_i+\lambda R_j$ ($i\neq j$). Prove that $S$ and $S'$ have the same solutions.
\end{minipage}\par\nopagebreak
\begin{tcolorbox}[breakable,colback=ans,colframe=ansb,boxrule=0.5pt,left=4pt,right=4pt,top=2pt,bottom=2pt]\begin{enumerate}[leftmargin=2.2em,itemsep=1pt,topsep=1pt,label=\textbf{Step \arabic*.}]\item Let the equations of $S$ be $E_1,\dots,E_m$ (each equation is an expression $E_k$ equal to a number $\beta_k$). $S'$ has the same equations except that $E_i$ is replaced by $E_i+\lambda E_j$ (and $\beta_i$ by $\beta_i+\lambda\beta_j$).
\item ($\Rightarrow$) Suppose $\mathbf x$ solves every equation of $S$. Then $E_i(\mathbf x)=\beta_i$ and $E_j(\mathbf x)=\beta_j$. Hence $(E_i+\lambda E_j)(\mathbf x)=\beta_i+\lambda\beta_j$, so $\mathbf x$ also satisfies the new equation. All other equations are unchanged, so $\mathbf x$ solves $S'$.
\item ($\Leftarrow$) Suppose $\mathbf x$ solves $S'$. Then it satisfies $E_j=\beta_j$ (unchanged) and $E_i+\lambda E_j=\beta_i+\lambda\beta_j$. Subtract $\lambda$ times the first from the second: $E_i(\mathbf x)=\beta_i$. So $\mathbf x$ solves $S$.
\item Both directions hold, so the solution sets are equal. The same argument (or the observation that each operation can be undone by another operation of the same kind) covers swaps and scaling by $c\neq0$.
\end{enumerate}\textbf{Answer:} Each solution of $S$ solves $S'$, and subtracting $\lambda E_j$ recovers $E_i$, so each solution of $S'$ solves $S$; the solution sets are equal.\end{tcolorbox}
\par\vspace{2pt}
\Needspace{16\baselineskip}\subsection*{\normalsize\color{ink} Part B --- Solving Linear Systems}
\Needspace{9\baselineskip}\begin{minipage}{\linewidth}\textbf{7.} \hfill{\scriptsize\color{ink}[Foundational]}\par\nopagebreak Solve by row reduction: $\begin{cases}x+y=3\\ 2x-y=0\end{cases}$
\end{minipage}\par\nopagebreak
\begin{tcolorbox}[breakable,colback=ans,colframe=ansb,boxrule=0.5pt,left=4pt,right=4pt,top=2pt,bottom=2pt]\begin{enumerate}[leftmargin=2.2em,itemsep=1pt,topsep=1pt,label=\textbf{Step \arabic*.}]\item Write the system as an augmented matrix, one row per equation, with a bar before the right-hand side:\par\smallskip $\displaystyle \left(\begin{array}{cc|c}1 & 1 & 3 \\ 2 & -1 & 0\end{array}\right)$
\item Apply $R_2\to R_2-2R_1$:\par\smallskip $\displaystyle \left(\begin{array}{cc|c}1 & 1 & 3 \\ 0 & -3 & -6\end{array}\right)$
\item Apply $R_2\to \dfrac{1}{-3} R_2$:\par\smallskip $\displaystyle \left(\begin{array}{cc|c}1 & 1 & 3 \\ 0 & 1 & 2\end{array}\right)$
\item Apply $R_1\to R_1-R_2$:\par\smallskip $\displaystyle \left(\begin{array}{cc|c}1 & 0 & 1 \\ 0 & 1 & 2\end{array}\right)$
\item Every column of the coefficient part holds a pivot, so there are no free variables and the solution is unique.
\item Read each nonzero row as an equation for its basic variable:\par\smallskip $\displaystyle\begin{aligned}x&=1\\ y&=2\end{aligned}$
\item Check by substituting $(x,y)=(1,2)$ into the original equations: the left sides give $(3,0)$, the right-hand sides \checkmark.
\end{enumerate}\textbf{Answer:} Unique solution: $(x,y)=(1,2)$.\end{tcolorbox}
\par\vspace{2pt}
\Needspace{9\baselineskip}\begin{minipage}{\linewidth}\textbf{8.} \hfill{\scriptsize\color{ink}[Foundational]}\par\nopagebreak Solve by row reduction: $\begin{cases}x+y+2z=5\\ 2x-y+z=4\\ x+2y-z=3\end{cases}$
\end{minipage}\par\nopagebreak
\begin{tcolorbox}[breakable,colback=ans,colframe=ansb,boxrule=0.5pt,left=4pt,right=4pt,top=2pt,bottom=2pt]\begin{enumerate}[leftmargin=2.2em,itemsep=1pt,topsep=1pt,label=\textbf{Step \arabic*.}]\item Write the system as an augmented matrix, one row per equation, with a bar before the right-hand side:\par\smallskip $\displaystyle \left(\begin{array}{ccc|c}1 & 1 & 2 & 5 \\ 2 & -1 & 1 & 4 \\ 1 & 2 & -1 & 3\end{array}\right)$
\item Apply $R_2\to R_2-2R_1,\ R_3\to R_3-R_1$:\par\smallskip $\displaystyle \left(\begin{array}{ccc|c}1 & 1 & 2 & 5 \\ 0 & -3 & -3 & -6 \\ 0 & 1 & -3 & -2\end{array}\right)$
\item Apply $R_2\to \dfrac{1}{-3} R_2$:\par\smallskip $\displaystyle \left(\begin{array}{ccc|c}1 & 1 & 2 & 5 \\ 0 & 1 & 1 & 2 \\ 0 & 1 & -3 & -2\end{array}\right)$
\item Apply $R_1\to R_1-R_2,\ R_3\to R_3-R_2$:\par\smallskip $\displaystyle \left(\begin{array}{ccc|c}1 & 0 & 1 & 3 \\ 0 & 1 & 1 & 2 \\ 0 & 0 & -4 & -4\end{array}\right)$
\item Apply $R_3\to \dfrac{1}{-4} R_3$:\par\smallskip $\displaystyle \left(\begin{array}{ccc|c}1 & 0 & 1 & 3 \\ 0 & 1 & 1 & 2 \\ 0 & 0 & 1 & 1\end{array}\right)$
\item Apply $R_1\to R_1-R_3,\ R_2\to R_2-R_3$:\par\smallskip $\displaystyle \left(\begin{array}{ccc|c}1 & 0 & 0 & 2 \\ 0 & 1 & 0 & 1 \\ 0 & 0 & 1 & 1\end{array}\right)$
\item Every column of the coefficient part holds a pivot, so there are no free variables and the solution is unique.
\item Read each nonzero row as an equation for its basic variable:\par\smallskip $\displaystyle\begin{aligned}x&=2\\ y&=1\\ z&=1\end{aligned}$
\item Check by substituting $(x,y,z)=(2,1,1)$ into the original equations: the left sides give $(5,4,3)$, the right-hand sides \checkmark.
\end{enumerate}\textbf{Answer:} Unique solution: $(x,y,z)=(2,1,1)$.\end{tcolorbox}
\par\vspace{2pt}
\Needspace{9\baselineskip}\begin{minipage}{\linewidth}\textbf{9.} \hfill{\scriptsize\color{ink}[Intermediate]}\par\nopagebreak Solve: $\begin{cases}x+2y-z=4\\ 2x+5y+z=9\end{cases}$ Describe the solution set in parametric vector form.
\end{minipage}\par\nopagebreak
\begin{tcolorbox}[breakable,colback=ans,colframe=ansb,boxrule=0.5pt,left=4pt,right=4pt,top=2pt,bottom=2pt]\begin{enumerate}[leftmargin=2.2em,itemsep=1pt,topsep=1pt,label=\textbf{Step \arabic*.}]\item Write the system as an augmented matrix, one row per equation, with a bar before the right-hand side:\par\smallskip $\displaystyle \left(\begin{array}{ccc|c}1 & 2 & -1 & 4 \\ 2 & 5 & 1 & 9\end{array}\right)$
\item Apply $R_2\to R_2-2R_1$:\par\smallskip $\displaystyle \left(\begin{array}{ccc|c}1 & 2 & -1 & 4 \\ 0 & 1 & 3 & 1\end{array}\right)$
\item Apply $R_1\to R_1-2R_2$:\par\smallskip $\displaystyle \left(\begin{array}{ccc|c}1 & 0 & -7 & 2 \\ 0 & 1 & 3 & 1\end{array}\right)$
\item Pivot columns are the columns of the basic variables: $x$, $y$. The remaining variable(s) $z$ are \emph{free}: call them $z=s$.
\item Read each nonzero row as an equation for its basic variable:\par\smallskip $\displaystyle\begin{aligned}x&=2+7s\\ y&=1-3s\\ z&=s\end{aligned}$
\item Write this as a vector: $\displaystyle \mathbf x=\begin{pmatrix}2\\1\\0\end{pmatrix}+s\begin{pmatrix}7\\-3\\1\end{pmatrix}$ (the parameters $s$ may be any numbers).
\item Check with $s=1$: $\mathbf x=\begin{pmatrix}9\\-2\\1\end{pmatrix}$, and substituting into the original equations gives the right-hand sides $(4,9)$ \checkmark.
\end{enumerate}\textbf{Answer:} Infinitely many solutions: $\mathbf x=\begin{pmatrix}2\\1\\0\end{pmatrix}+s\begin{pmatrix}7\\-3\\1\end{pmatrix}$.\end{tcolorbox}
\par\vspace{2pt}
\Needspace{9\baselineskip}\begin{minipage}{\linewidth}\textbf{10.} \hfill{\scriptsize\color{ink}[Intermediate]}\par\nopagebreak Decide whether the system is consistent: $\begin{cases}x+2y+z=1\\ 2x+4y+3z=3\\ x+2y-z=2\end{cases}$
\end{minipage}\par\nopagebreak
\begin{tcolorbox}[breakable,colback=ans,colframe=ansb,boxrule=0.5pt,left=4pt,right=4pt,top=2pt,bottom=2pt]\begin{enumerate}[leftmargin=2.2em,itemsep=1pt,topsep=1pt,label=\textbf{Step \arabic*.}]\item Write the system as an augmented matrix, one row per equation, with a bar before the right-hand side:\par\smallskip $\displaystyle \left(\begin{array}{ccc|c}1 & 2 & 1 & 1 \\ 2 & 4 & 3 & 3 \\ 1 & 2 & -1 & 2\end{array}\right)$
\item Apply $R_2\to R_2-2R_1,\ R_3\to R_3-R_1$:\par\smallskip $\displaystyle \left(\begin{array}{ccc|c}1 & 2 & 1 & 1 \\ 0 & 0 & 1 & 1 \\ 0 & 0 & -2 & 1\end{array}\right)$
\item Apply $R_1\to R_1-R_2,\ R_3\to R_3+2R_2$:\par\smallskip $\displaystyle \left(\begin{array}{ccc|c}1 & 2 & 0 & 0 \\ 0 & 0 & 1 & 1 \\ 0 & 0 & 0 & 3\end{array}\right)$
\item Row 3 reads $0=1$, a contradiction. No numbers $x,y,z$ can satisfy it, so the system is \emph{inconsistent}.
\end{enumerate}\textbf{Answer:} No solution (inconsistent).\end{tcolorbox}
\par\vspace{2pt}
\Needspace{9\baselineskip}\begin{minipage}{\linewidth}\textbf{11.} \hfill{\scriptsize\color{ink}[Intermediate]}\par\nopagebreak Solve $\begin{cases}x+y+z+w=2\\ x+2y+3w=5\end{cases}$ and identify the basic and free variables.
\end{minipage}\par\nopagebreak
\begin{tcolorbox}[breakable,colback=ans,colframe=ansb,boxrule=0.5pt,left=4pt,right=4pt,top=2pt,bottom=2pt]\begin{enumerate}[leftmargin=2.2em,itemsep=1pt,topsep=1pt,label=\textbf{Step \arabic*.}]\item Write the system as an augmented matrix, one row per equation, with a bar before the right-hand side:\par\smallskip $\displaystyle \left(\begin{array}{cccc|c}1 & 1 & 1 & 1 & 2 \\ 1 & 2 & 0 & 3 & 5\end{array}\right)$
\item Apply $R_2\to R_2-R_1$:\par\smallskip $\displaystyle \left(\begin{array}{cccc|c}1 & 1 & 1 & 1 & 2 \\ 0 & 1 & -1 & 2 & 3\end{array}\right)$
\item Apply $R_1\to R_1-R_2$:\par\smallskip $\displaystyle \left(\begin{array}{cccc|c}1 & 0 & 2 & -1 & -1 \\ 0 & 1 & -1 & 2 & 3\end{array}\right)$
\item Pivot columns are the columns of the basic variables: $x$, $y$. The remaining variable(s) $z$, $w$ are \emph{free}: call them $z=s$, $w=t$.
\item Read each nonzero row as an equation for its basic variable:\par\smallskip $\displaystyle\begin{aligned}x&=-1-2s+t\\ y&=3+s-2t\\ z&=s\\ w&=t\end{aligned}$
\item Write this as a vector: $\displaystyle \mathbf x=\begin{pmatrix}-1\\3\\0\\0\end{pmatrix}+s\begin{pmatrix}-2\\1\\1\\0\end{pmatrix}+t\begin{pmatrix}1\\-2\\0\\1\end{pmatrix}$ (the parameters $s$, $t$ may be any numbers).
\item Check with $s=1$, $t=1$: $\mathbf x=\begin{pmatrix}-2\\2\\1\\1\end{pmatrix}$, and substituting into the original equations gives the right-hand sides $(2,5)$ \checkmark.
\end{enumerate}\textbf{Answer:} Infinitely many solutions: $\mathbf x=\begin{pmatrix}-1\\3\\0\\0\end{pmatrix}+s\begin{pmatrix}-2\\1\\1\\0\end{pmatrix}+t\begin{pmatrix}1\\-2\\0\\1\end{pmatrix}$.\end{tcolorbox}
\par\vspace{2pt}
\Needspace{9\baselineskip}\begin{minipage}{\linewidth}\textbf{12.} \hfill{\scriptsize\color{ink}[Intermediate]}\par\nopagebreak Solve the homogeneous system $\begin{cases}x+2y-w=0\\ 2x+4y+z+w=0\\ x+2y+z+2w=0\end{cases}$ and write the solution set as a span.
\end{minipage}\par\nopagebreak
\begin{tcolorbox}[breakable,colback=ans,colframe=ansb,boxrule=0.5pt,left=4pt,right=4pt,top=2pt,bottom=2pt]\begin{enumerate}[leftmargin=2.2em,itemsep=1pt,topsep=1pt,label=\textbf{Step \arabic*.}]\item Write the system as an augmented matrix, one row per equation, with a bar before the right-hand side:\par\smallskip $\displaystyle \left(\begin{array}{cccc|c}1 & 2 & 0 & -1 & 0 \\ 2 & 4 & 1 & 1 & 0 \\ 1 & 2 & 1 & 2 & 0\end{array}\right)$
\item Apply $R_2\to R_2-2R_1,\ R_3\to R_3-R_1$:\par\smallskip $\displaystyle \left(\begin{array}{cccc|c}1 & 2 & 0 & -1 & 0 \\ 0 & 0 & 1 & 3 & 0 \\ 0 & 0 & 1 & 3 & 0\end{array}\right)$
\item Apply $R_3\to R_3-R_2$:\par\smallskip $\displaystyle \left(\begin{array}{cccc|c}1 & 2 & 0 & -1 & 0 \\ 0 & 0 & 1 & 3 & 0 \\ 0 & 0 & 0 & 0 & 0\end{array}\right)$
\item Pivot columns are the columns of the basic variables: $x$, $z$. The remaining variable(s) $y$, $w$ are \emph{free}: call them $y=s$, $w=t$.
\item Read each nonzero row as an equation for its basic variable:\par\smallskip $\displaystyle\begin{aligned}x&=-2s+t\\ y&=s\\ z&=-3t\\ w&=t\end{aligned}$
\item Write this as a vector: $\displaystyle \mathbf x=s\begin{pmatrix}-2\\1\\0\\0\end{pmatrix}+t\begin{pmatrix}1\\0\\-3\\1\end{pmatrix}$ (the parameters $s$, $t$ may be any numbers).
\item Check with $s=1$, $t=1$: $\mathbf x=\begin{pmatrix}-1\\1\\-3\\1\end{pmatrix}$, and substituting into the original equations gives the right-hand sides $(0,0,0)$ \checkmark.
\end{enumerate}\textbf{Answer:} Infinitely many solutions: $\mathbf x=s\begin{pmatrix}-2\\1\\0\\0\end{pmatrix}+t\begin{pmatrix}1\\0\\-3\\1\end{pmatrix}$. Equivalently the solution set is the span of the vectors multiplying the parameters (a subspace of $\mathbb{R}^{4}$; it always contains $\mathbf 0$).\end{tcolorbox}
\par\vspace{2pt}
\Needspace{9\baselineskip}\begin{minipage}{\linewidth}\textbf{13.} \hfill{\scriptsize\color{ink}[Intermediate]}\par\nopagebreak Over $\mathbb{C}$, solve $\begin{cases}x+iy=1+i\\ ix+2y=i\end{cases}$
\end{minipage}\par\nopagebreak
\begin{tcolorbox}[breakable,colback=ans,colframe=ansb,boxrule=0.5pt,left=4pt,right=4pt,top=2pt,bottom=2pt]\begin{enumerate}[leftmargin=2.2em,itemsep=1pt,topsep=1pt,label=\textbf{Step \arabic*.}]\item Write the system as an augmented matrix, one row per equation, with a bar before the right-hand side:\par\smallskip $\displaystyle \left(\begin{array}{cc|c}1 & i & 1+i \\ i & 2 & i\end{array}\right)$
\item Apply $R_2\to R_2-iR_1$:\par\smallskip $\displaystyle \left(\begin{array}{cc|c}1 & i & 1+i \\ 0 & 3 & 1\end{array}\right)$
\item Apply $R_2\to \dfrac{1}{3} R_2$:\par\smallskip $\displaystyle \left(\begin{array}{cc|c}1 & i & 1+i \\ 0 & 1 & \dfrac{1}{3}\end{array}\right)$
\item Apply $R_1\to R_1-iR_2$:\par\smallskip $\displaystyle \left(\begin{array}{cc|c}1 & 0 & 1+\dfrac{2}{3}i \\ 0 & 1 & \dfrac{1}{3}\end{array}\right)$
\item Every column of the coefficient part holds a pivot, so there are no free variables and the solution is unique.
\item Read each nonzero row as an equation for its basic variable:\par\smallskip $\displaystyle\begin{aligned}x&=1+\dfrac{2}{3}i\\ y&=\dfrac{1}{3}\end{aligned}$
\item Check by substituting $(x,y)=(1+\dfrac{2}{3}i,\dfrac{1}{3})$ into the original equations: the left sides give $(1+i,i)$, the right-hand sides \checkmark.
\end{enumerate}\textbf{Answer:} Unique solution: $(x,y)=(1+\dfrac{2}{3}i,\dfrac{1}{3})$.\end{tcolorbox}
\par\vspace{2pt}
\Needspace{9\baselineskip}\begin{minipage}{\linewidth}\textbf{14.} \hfill{\scriptsize\color{ink}[Challenge]}\par\nopagebreak For which real $k$ does the system $\begin{cases}x+y+z=1\\ x+2y+3z=2\\ x+3y+kz=4\end{cases}$ have (i) a unique solution, (ii) no solution, (iii) infinitely many solutions?
\end{minipage}\par\nopagebreak
\begin{tcolorbox}[breakable,colback=ans,colframe=ansb,boxrule=0.5pt,left=4pt,right=4pt,top=2pt,bottom=2pt]\begin{enumerate}[leftmargin=2.2em,itemsep=1pt,topsep=1pt,label=\textbf{Step \arabic*.}]\item Write the augmented matrix and reduce, keeping $k$ as a symbol: $\displaystyle\left(\begin{array}{ccc|c}1&1&1&1\\1&2&3&2\\1&3&k&4\end{array}\right)$.
\item Apply $R_2\to R_2-R_1,\ R_3\to R_3-R_1$: $\displaystyle\left(\begin{array}{ccc|c}1&1&1&1\\0&1&2&1\\0&2&k-1&3\end{array}\right)$.
\item Apply $R_3\to R_3-2R_2$: $\displaystyle\left(\begin{array}{ccc|c}1&1&1&1\\0&1&2&1\\0&0&k-5&1\end{array}\right)$.
\item Read the last row: $(k-5)z=1$. If $k\neq5$ we can divide by $k-5$, so $z=\dfrac{1}{k-5}$; then $y$ and $x$ follow by back-substitution. There are three pivots, so the solution is unique.
\item If $k=5$ the last row reads $0=1$, a contradiction: no solution.
\item Infinitely many solutions would need a free variable and a consistent system. For $k=5$ the system is inconsistent, and for $k\neq5$ there is no free variable. So this never happens.
\item Check one value: $k=6$ gives $z=1$, then $y+2=1$ so $y=-1$, then $x-1+1=1$ so $x=1$; test in equation 3: $1+3(-1)+6(1)=4$ \checkmark.
\end{enumerate}\textbf{Answer:} (i) unique solution for every $k\neq5$; (ii) no solution when $k=5$; (iii) never infinitely many.\end{tcolorbox}
\par\vspace{2pt}
\Needspace{9\baselineskip}\begin{minipage}{\linewidth}\textbf{15.} \hfill{\scriptsize\color{ink}[Challenge]}\par\nopagebreak A homogeneous system has $m$ equations and $n$ unknowns with $n>m$. Prove it has a nonzero solution.
\end{minipage}\par\nopagebreak
\begin{tcolorbox}[breakable,colback=ans,colframe=ansb,boxrule=0.5pt,left=4pt,right=4pt,top=2pt,bottom=2pt]\begin{enumerate}[leftmargin=2.2em,itemsep=1pt,topsep=1pt,label=\textbf{Step \arabic*.}]\item Reduce the $m\times n$ coefficient matrix to RREF. Each pivot occupies its own row and its own column, so there are at most $m$ pivots (one per row).
\item The coefficient matrix has $n$ columns and $n>m\ge$ (number of pivots). So at least $n-m\ge1$ columns contain no pivot; the variables in those columns are free.
\item A homogeneous system is always consistent (the zero vector solves it). Give one free variable the value $1$ (others $0$) and solve for the basic variables from the RREF rows. The resulting solution has a coordinate equal to $1$, so it is nonzero.
\end{enumerate}\textbf{Answer:} At most $m$ pivot columns, so at least $n-m\ge1$ free variables; setting a free variable to $1$ gives a nonzero solution of the always-consistent homogeneous system.\end{tcolorbox}
\par\vspace{2pt}
\Needspace{16\baselineskip}\subsection*{\normalsize\color{ink} Part C --- Column Spaces, Relations and Proofs}
\Needspace{9\baselineskip}\begin{minipage}{\linewidth}\textbf{16.} \hfill{\scriptsize\color{ink}[Foundational]}\par\nopagebreak The matrix $R=\begin{pmatrix}1 & 0 & 3 & 0 & -2 \\ 0 & 1 & -1 & 0 & 4 \\ 0 & 0 & 0 & 1 & 5\end{pmatrix}$ is the RREF of the coefficient matrix of a homogeneous system. State the pivot columns, the basic and free variables (use $x_1,\dots,x_5$), and the general solution of $R\mathbf x=\mathbf 0$.
\end{minipage}\par\nopagebreak
\begin{tcolorbox}[breakable,colback=ans,colframe=ansb,boxrule=0.5pt,left=4pt,right=4pt,top=2pt,bottom=2pt]\begin{enumerate}[leftmargin=2.2em,itemsep=1pt,topsep=1pt,label=\textbf{Step \arabic*.}]\item Pivot columns are the columns containing a leading $1$: columns $1$, $2$ and $4$. So $x_1,x_2,x_4$ are basic and $x_3,x_5$ are free; set $x_3=s$, $x_5=t$.
\item Read the rows: row 1: $x_1+3x_3-2x_5=0$, so $x_1=-3s+2t$. Row 2: $x_2-x_3+4x_5=0$, so $x_2=s-4t$. Row 3: $x_4+5x_5=0$, so $x_4=-5t$.
\item Vector form: $\mathbf x=\begin{pmatrix}-3s+2t\\ s-4t\\ s\\ -5t\\ t\end{pmatrix}=s\begin{pmatrix}-3\\1\\1\\0\\0\end{pmatrix}+t\begin{pmatrix}2\\-4\\0\\-5\\1\end{pmatrix}$.
\item Check the first vector in row 1: $(-3)+3(1)-2(0)=0$ \checkmark; in row 2: $1-1+0=0$ \checkmark; in row 3: $0+0=0$ \checkmark.
\end{enumerate}\textbf{Answer:} Pivots: columns $1,2,4$. Free: $x_3,x_5$. $\mathbf x=s(-3,1,1,0,0)^{T}+t(2,-4,0,-5,1)^{T}$.\end{tcolorbox}
\par\vspace{2pt}
\Needspace{9\baselineskip}\begin{minipage}{\linewidth}\textbf{17.} \hfill{\scriptsize\color{ink}[Intermediate]}\par\nopagebreak Let $A=\begin{pmatrix}1 & 2 & 0 & 1 \\ 2 & 4 & 1 & 3 \\ 1 & 2 & 1 & 2\end{pmatrix}$. (a) Reduce $A$ to RREF. (b) Give a basis of the column space $\operatorname{col}(A)$ made of columns of $A$ itself. (c) Express each non-pivot column of $A$ as a linear combination of the pivot columns.
\end{minipage}\par\nopagebreak
\begin{tcolorbox}[breakable,colback=ans,colframe=ansb,boxrule=0.5pt,left=4pt,right=4pt,top=2pt,bottom=2pt]\begin{enumerate}[leftmargin=2.2em,itemsep=1pt,topsep=1pt,label=\textbf{Step \arabic*.}]\item (a) Row reduce.
\item Apply $R_2\to R_2-2R_1,\ R_3\to R_3-R_1$:\par\smallskip $\displaystyle \begin{pmatrix}1 & 2 & 0 & 1 \\ 0 & 0 & 1 & 1 \\ 0 & 0 & 1 & 1\end{pmatrix}$
\item Apply $R_3\to R_3-R_2$:\par\smallskip $\displaystyle \begin{pmatrix}1 & 2 & 0 & 1 \\ 0 & 0 & 1 & 1 \\ 0 & 0 & 0 & 0\end{pmatrix}$
\item The RREF has pivots in columns $1$ and $3$.
\item (b) Rule: the pivot columns of the \emph{original} matrix $A$ (not of the RREF) form a basis of $\operatorname{col}(A)$. So a basis is $\left\{\begin{pmatrix}1\\2\\1\end{pmatrix},\ \begin{pmatrix}0\\1\\1\end{pmatrix}\right\}$.
\item (c) The entries of a non-pivot column of the RREF are the coefficients: column 2 of the RREF is $\begin{pmatrix}2\\0\end{pmatrix}$, so $\text{col}_2=2\,\text{col}_1$; column 4 of the RREF is $\begin{pmatrix}1\\1\end{pmatrix}$, so $\text{col}_4=1\,\text{col}_1+1\,\text{col}_3$.
\item Check column 4: $1\cdot\begin{pmatrix}1\\2\\1\end{pmatrix}+1\cdot\begin{pmatrix}0\\1\\1\end{pmatrix}=\begin{pmatrix}1\\3\\2\end{pmatrix}$ \checkmark.
\end{enumerate}\textbf{Answer:} RREF $=\begin{pmatrix}1 & 2 & 0 & 1 \\ 0 & 0 & 1 & 1 \\ 0 & 0 & 0 & 0\end{pmatrix}$. Basis of $\operatorname{col}(A)$: $\{\begin{pmatrix}1\\2\\1\end{pmatrix},\begin{pmatrix}0\\1\\1\end{pmatrix}\}$ (columns $1$ and $3$ of $A$); $\text{col}_2=2\,\text{col}_1$, $\text{col}_4=\text{col}_1+\text{col}_3$.\end{tcolorbox}
\par\vspace{2pt}
\Needspace{9\baselineskip}\begin{minipage}{\linewidth}\textbf{18.} \hfill{\scriptsize\color{ink}[Intermediate]}\par\nopagebreak For $A=\begin{pmatrix}1 & 2 & 3 \\ 2 & 4 & 7 \\ 3 & 6 & 10\end{pmatrix}$, reduce to RREF and give a basis of the row space $\operatorname{row}(A)$ and the number of pivots.
\end{minipage}\par\nopagebreak
\begin{tcolorbox}[breakable,colback=ans,colframe=ansb,boxrule=0.5pt,left=4pt,right=4pt,top=2pt,bottom=2pt]\begin{enumerate}[leftmargin=2.2em,itemsep=1pt,topsep=1pt,label=\textbf{Step \arabic*.}]\item Row operations replace each row by a combination of rows, and they can be reversed, so they do not change the span of the rows. Hence $\operatorname{row}(A)=\operatorname{row}(R)$ where $R$ is the RREF.
\item Apply $R_2\to R_2-2R_1,\ R_3\to R_3-3R_1$:\par\smallskip $\displaystyle \begin{pmatrix}1 & 2 & 3 \\ 0 & 0 & 1 \\ 0 & 0 & 1\end{pmatrix}$
\item Apply $R_1\to R_1-3R_2,\ R_3\to R_3-R_2$:\par\smallskip $\displaystyle \begin{pmatrix}1 & 2 & 0 \\ 0 & 0 & 1 \\ 0 & 0 & 0\end{pmatrix}$
\item The nonzero rows of the RREF are linearly independent (each has a $1$ in a pivot column where all the others have $0$) and span the row space. So they form a basis.
\end{enumerate}\textbf{Answer:} RREF $=\begin{pmatrix}1 & 2 & 0 \\ 0 & 0 & 1 \\ 0 & 0 & 0\end{pmatrix}$; basis of $\operatorname{row}(A)$: $(1,2,0)$, $(0,0,1)$; the rank (number of pivots) is $2$.\end{tcolorbox}
\par\vspace{2pt}
\Needspace{9\baselineskip}\begin{minipage}{\linewidth}\textbf{19.} \hfill{\scriptsize\color{ink}[Intermediate]}\par\nopagebreak Decide whether $\mathbf w=(3,2,1)$ lies in $\operatorname{span}\{(1,0,1),(2,1,0),(0,1,-2)\}\subseteq\mathbb{R}^{3}$ by row reducing one augmented matrix. Give the coefficients if it does.
\end{minipage}\par\nopagebreak
\begin{tcolorbox}[breakable,colback=ans,colframe=ansb,boxrule=0.5pt,left=4pt,right=4pt,top=2pt,bottom=2pt]\begin{enumerate}[leftmargin=2.2em,itemsep=1pt,topsep=1pt,label=\textbf{Step \arabic*.}]\item Write the system as an augmented matrix, one row per equation, with a bar before the right-hand side:\par\smallskip $\displaystyle \left(\begin{array}{ccc|c}1 & 2 & 0 & 3 \\ 0 & 1 & 1 & 2 \\ 1 & 0 & -2 & 1\end{array}\right)$
\item (The columns are the three vectors; the right-hand side is $\mathbf w$.)
\item Apply $R_3\to R_3-R_1$:\par\smallskip $\displaystyle \left(\begin{array}{ccc|c}1 & 2 & 0 & 3 \\ 0 & 1 & 1 & 2 \\ 0 & -2 & -2 & -2\end{array}\right)$
\item Apply $R_1\to R_1-2R_2,\ R_3\to R_3+2R_2$:\par\smallskip $\displaystyle \left(\begin{array}{ccc|c}1 & 0 & -2 & -1 \\ 0 & 1 & 1 & 2 \\ 0 & 0 & 0 & 2\end{array}\right)$
\item Row 3 reads $0=1$, a contradiction. No numbers $a,b,c$ can satisfy it, so the system is \emph{inconsistent}.
\end{enumerate}\textbf{Answer:} No: the system is inconsistent, so $\mathbf w\notin\operatorname{span}$.\end{tcolorbox}
\par\vspace{2pt}
\Needspace{9\baselineskip}\begin{minipage}{\linewidth}\textbf{20.} \hfill{\scriptsize\color{ink}[Intermediate]}\par\nopagebreak Place the vectors $(1,2,1,0),\ (1,0,1,1),\ (2,1,0,1)$ of $\mathbb{R}^{4}$ as columns of a matrix. Use RREF to decide whether they are linearly independent.
\end{minipage}\par\nopagebreak
\begin{tcolorbox}[breakable,colback=ans,colframe=ansb,boxrule=0.5pt,left=4pt,right=4pt,top=2pt,bottom=2pt]\begin{enumerate}[leftmargin=2.2em,itemsep=1pt,topsep=1pt,label=\textbf{Step \arabic*.}]\item Put the vectors in the columns of $A$; independence means $A\mathbf a=\mathbf 0$ has only the zero solution.
\item $A=\begin{pmatrix}1 & 1 & 2 \\ 2 & 0 & 1 \\ 1 & 1 & 0 \\ 0 & 1 & 1\end{pmatrix}$. Reduce it.
\item Apply $R_2\to R_2-2R_1,\ R_3\to R_3-R_1$:\par\smallskip $\displaystyle \begin{pmatrix}1 & 1 & 2 \\ 0 & -2 & -3 \\ 0 & 0 & -2 \\ 0 & 1 & 1\end{pmatrix}$
\item Apply $R_2\to \dfrac{1}{-2} R_2$:\par\smallskip $\displaystyle \begin{pmatrix}1 & 1 & 2 \\ 0 & 1 & \dfrac{3}{2} \\ 0 & 0 & -2 \\ 0 & 1 & 1\end{pmatrix}$
\item Apply $R_1\to R_1-R_2,\ R_4\to R_4-R_2$:\par\smallskip $\displaystyle \begin{pmatrix}1 & 0 & \dfrac{1}{2} \\ 0 & 1 & \dfrac{3}{2} \\ 0 & 0 & -2 \\ 0 & 0 & -\dfrac{1}{2}\end{pmatrix}$
\item Apply $R_3\to \dfrac{1}{-2} R_3$:\par\smallskip $\displaystyle \begin{pmatrix}1 & 0 & \dfrac{1}{2} \\ 0 & 1 & \dfrac{3}{2} \\ 0 & 0 & 1 \\ 0 & 0 & -\dfrac{1}{2}\end{pmatrix}$
\item Apply $R_1\to R_1-\dfrac{1}{2}R_3,\ R_2\to R_2-\dfrac{3}{2}R_3,\ R_4\to R_4+\dfrac{1}{2}R_3$:\par\smallskip $\displaystyle \begin{pmatrix}1 & 0 & 0 \\ 0 & 1 & 0 \\ 0 & 0 & 1 \\ 0 & 0 & 0\end{pmatrix}$
\item Every column holds a pivot, so there are no free variables: the only solution of $A\mathbf a=\mathbf 0$ is $\mathbf a=\mathbf 0$.
\end{enumerate}\textbf{Answer:} Linearly independent (a pivot in every column).\end{tcolorbox}
\par\vspace{2pt}
\Needspace{9\baselineskip}\begin{minipage}{\linewidth}\textbf{21.} \hfill{\scriptsize\color{ink}[Intermediate]}\par\nopagebreak Show that $(1,0,1),(1,1,0),(0,1,1),(2,3,1)$ in $\mathbb{R}^{3}$ are linearly dependent, and use RREF to find an explicit relation among them.
\end{minipage}\par\nopagebreak
\begin{tcolorbox}[breakable,colback=ans,colframe=ansb,boxrule=0.5pt,left=4pt,right=4pt,top=2pt,bottom=2pt]\begin{enumerate}[leftmargin=2.2em,itemsep=1pt,topsep=1pt,label=\textbf{Step \arabic*.}]\item Four vectors in $\mathbb{R}^{3}$: the matrix with these columns is $3\times4$, so it cannot have a pivot in every column. Dependence is guaranteed; we now find a relation.
\item $A=\begin{pmatrix}1 & 1 & 0 & 2 \\ 0 & 1 & 1 & 3 \\ 1 & 0 & 1 & 1\end{pmatrix}$. Reduce it.
\item Apply $R_3\to R_3-R_1$:\par\smallskip $\displaystyle \begin{pmatrix}1 & 1 & 0 & 2 \\ 0 & 1 & 1 & 3 \\ 0 & -1 & 1 & -1\end{pmatrix}$
\item Apply $R_1\to R_1-R_2,\ R_3\to R_3+R_2$:\par\smallskip $\displaystyle \begin{pmatrix}1 & 0 & -1 & -1 \\ 0 & 1 & 1 & 3 \\ 0 & 0 & 2 & 2\end{pmatrix}$
\item Apply $R_3\to \dfrac{1}{2} R_3$:\par\smallskip $\displaystyle \begin{pmatrix}1 & 0 & -1 & -1 \\ 0 & 1 & 1 & 3 \\ 0 & 0 & 1 & 1\end{pmatrix}$
\item Apply $R_1\to R_1+R_3,\ R_2\to R_2-R_3$:\par\smallskip $\displaystyle \begin{pmatrix}1 & 0 & 0 & 0 \\ 0 & 1 & 0 & 2 \\ 0 & 0 & 1 & 1\end{pmatrix}$
\item Column 4 has no pivot, so $a_4$ is free: set $a_4=1$. Solving the RREF rows gives the vector $\begin{pmatrix}0\\-2\\-1\\1\end{pmatrix}$.
\item So the relation is $(0)v_1+(-2)v_2+(-1)v_3+(1)v_4=\mathbf 0$. Check: $(0)\begin{pmatrix}1\\0\\1\end{pmatrix}+(-2)\begin{pmatrix}1\\1\\0\end{pmatrix}+(-1)\begin{pmatrix}0\\1\\1\end{pmatrix}+(1)\begin{pmatrix}2\\3\\1\end{pmatrix}=\begin{pmatrix}0\\0\\0\end{pmatrix}$ \checkmark.
\end{enumerate}\textbf{Answer:} Dependent: $(0)v_1+(-2)v_2+(-1)v_3+(1)v_4=\mathbf 0$.\end{tcolorbox}
\par\vspace{2pt}
\Needspace{9\baselineskip}\begin{minipage}{\linewidth}\textbf{22.} \hfill{\scriptsize\color{ink}[Intermediate]}\par\nopagebreak In $P_{2}(\mathbb{R})$, decide whether $q(x)=2+3x+x^{2}\in\operatorname{span}\{1+x,\ x+x^{2},\ 1+x^{2}\}$. Set up the system, reduce, and conclude.
\end{minipage}\par\nopagebreak
\begin{tcolorbox}[breakable,colback=ans,colframe=ansb,boxrule=0.5pt,left=4pt,right=4pt,top=2pt,bottom=2pt]\begin{enumerate}[leftmargin=2.2em,itemsep=1pt,topsep=1pt,label=\textbf{Step \arabic*.}]\item We need $a,b,c$ with $a(1+x)+b(x+x^{2})+c(1+x^{2})=2+3x+x^{2}$. Expand and collect powers: $(a+c)+(a+b)x+(b+c)x^{2}$.
\item A polynomial equality means equal coefficients: constant $a+c=2$; $x$: $a+b=3$; $x^{2}$: $b+c=1$. In matrix form (columns are the coordinate vectors with respect to $1,x,x^{2}$).
\item Write the system as an augmented matrix, one row per equation, with a bar before the right-hand side:\par\smallskip $\displaystyle \left(\begin{array}{ccc|c}1 & 0 & 1 & 2 \\ 1 & 1 & 0 & 3 \\ 0 & 1 & 1 & 1\end{array}\right)$
\item Apply $R_2\to R_2-R_1$:\par\smallskip $\displaystyle \left(\begin{array}{ccc|c}1 & 0 & 1 & 2 \\ 0 & 1 & -1 & 1 \\ 0 & 1 & 1 & 1\end{array}\right)$
\item Apply $R_3\to R_3-R_2$:\par\smallskip $\displaystyle \left(\begin{array}{ccc|c}1 & 0 & 1 & 2 \\ 0 & 1 & -1 & 1 \\ 0 & 0 & 2 & 0\end{array}\right)$
\item Apply $R_3\to \dfrac{1}{2} R_3$:\par\smallskip $\displaystyle \left(\begin{array}{ccc|c}1 & 0 & 1 & 2 \\ 0 & 1 & -1 & 1 \\ 0 & 0 & 1 & 0\end{array}\right)$
\item Apply $R_1\to R_1-R_3,\ R_2\to R_2+R_3$:\par\smallskip $\displaystyle \left(\begin{array}{ccc|c}1 & 0 & 0 & 2 \\ 0 & 1 & 0 & 1 \\ 0 & 0 & 1 & 0\end{array}\right)$
\item Every column of the coefficient part holds a pivot, so there are no free variables and the solution is unique.
\item Read each nonzero row as an equation for its basic variable:\par\smallskip $\displaystyle\begin{aligned}a&=2\\ b&=1\\ c&=0\end{aligned}$
\item Check by substituting $(a,b,c)=(2,1,0)$ into the original equations: the left sides give $(2,3,1)$, the right-hand sides \checkmark.
\end{enumerate}\textbf{Answer:} Yes, $q$ is in the span; the coefficients are: $(a,b,c)=(2,1,0)$. That is, $q=2(1+x)+1(x+x^{2})+0(1+x^{2})$.\end{tcolorbox}
\par\vspace{2pt}
\Needspace{9\baselineskip}\begin{minipage}{\linewidth}\textbf{23.} \hfill{\scriptsize\color{ink}[Challenge]}\par\nopagebreak Let $A\mathbf x=\mathbf b$ be consistent and let $\mathbf p$ be one solution. Prove that the solution set is exactly $\{\mathbf p+\mathbf h:\ A\mathbf h=\mathbf 0\}$.
\end{minipage}\par\nopagebreak
\begin{tcolorbox}[breakable,colback=ans,colframe=ansb,boxrule=0.5pt,left=4pt,right=4pt,top=2pt,bottom=2pt]\begin{enumerate}[leftmargin=2.2em,itemsep=1pt,topsep=1pt,label=\textbf{Step \arabic*.}]\item Two inclusions. ($\supseteq$) Let $A\mathbf h=\mathbf 0$. Then $A(\mathbf p+\mathbf h)=A\mathbf p+A\mathbf h=\mathbf b+\mathbf 0=\mathbf b$, so $\mathbf p+\mathbf h$ is a solution.
\item ($\subseteq$) Let $\mathbf x$ be any solution of $A\mathbf x=\mathbf b$. Put $\mathbf h=\mathbf x-\mathbf p$. Then $A\mathbf h=A\mathbf x-A\mathbf p=\mathbf b-\mathbf b=\mathbf 0$, and $\mathbf x=\mathbf p+\mathbf h$.
\item Both inclusions hold, so the sets are equal. (This is why the general solution always has the shape ``particular solution $+$ solution of the homogeneous system''.)
\end{enumerate}\textbf{Answer:} $A(\mathbf p+\mathbf h)=\mathbf b$ for each $\mathbf h$ in the null set, and any solution $\mathbf x$ equals $\mathbf p+(\mathbf x-\mathbf p)$ with $A(\mathbf x-\mathbf p)=\mathbf 0$.\end{tcolorbox}
\par\vspace{2pt}
\Needspace{9\baselineskip}\begin{minipage}{\linewidth}\textbf{24.} \hfill{\scriptsize\color{ink}[Challenge]}\par\nopagebreak Find the condition on $a,b,c$ for the system $\begin{cases}x+2y-z=a\\ 2x+3y+z=b\\ x+y+2z=c\end{cases}$ to be consistent.
\end{minipage}\par\nopagebreak
\begin{tcolorbox}[breakable,colback=ans,colframe=ansb,boxrule=0.5pt,left=4pt,right=4pt,top=2pt,bottom=2pt]\begin{enumerate}[leftmargin=2.2em,itemsep=1pt,topsep=1pt,label=\textbf{Step \arabic*.}]\item Row reduce the augmented matrix, carrying $a,b,c$ along: $\displaystyle\left(\begin{array}{ccc|c}1&2&-1&a\\2&3&1&b\\1&1&2&c\end{array}\right)$.
\item $R_2\to R_2-2R_1,\ R_3\to R_3-R_1$: $\displaystyle\left(\begin{array}{ccc|c}1&2&-1&a\\0&-1&3&b-2a\\0&-1&3&c-a\end{array}\right)$.
\item $R_3\to R_3-R_2$: $\displaystyle\left(\begin{array}{ccc|c}1&2&-1&a\\0&-1&3&b-2a\\0&0&0&(c-a)-(b-2a)\end{array}\right)$. The last right-hand side simplifies to $a-b+c$.
\item The last row says $0=a-b+c$. If $a-b+c\neq0$ this is a contradiction (inconsistent). If $a-b+c=0$ the row is $0=0$ and the system is consistent, with a free variable (infinitely many solutions).
\item Check: take $(a,b,c)=(1,3,2)$, which satisfies $1-3+2=0$. Then $x=1-2y+z$ from equation 1; with $z=0$: $x=1-2y$, equation 2: $2-4y+3y=3$, so $y=-1$, $x=3$; equation 3: $3-1+0=2$ \checkmark.
\end{enumerate}\textbf{Answer:} Consistent if and only if $a-b+c=0$ (then infinitely many solutions); otherwise no solution.\end{tcolorbox}
\par\vspace{2pt}
\Needspace{9\baselineskip}\begin{minipage}{\linewidth}\textbf{25.} \hfill{\scriptsize\color{ink}[Challenge]}\par\nopagebreak Let $R$ be the RREF of $A$. Prove that every non-pivot column of $R$ is a linear combination of the pivot columns of $R$, and explain why the same relation holds among the corresponding columns of $A$.
\end{minipage}\par\nopagebreak
\begin{tcolorbox}[breakable,colback=ans,colframe=ansb,boxrule=0.5pt,left=4pt,right=4pt,top=2pt,bottom=2pt]\begin{enumerate}[leftmargin=2.2em,itemsep=1pt,topsep=1pt,label=\textbf{Step \arabic*.}]\item Let the pivot columns of $R$ be $j_1<\dots<j_r$. Column $j_k$ of $R$ is the standard vector $\mathbf e_k$ (a $1$ in row $k$, zeros elsewhere), because each pivot is $1$ and is the only nonzero entry in its column.
\item Any column $\mathbf c$ of $R$ has nonzero entries only in rows $1,\dots,r$ (rows below the last pivot are zero rows). So $\mathbf c=c_1\mathbf e_1+\dots+c_r\mathbf e_r=c_1R_{j_1}+\dots+c_rR_{j_r}$, where $c_k$ is the $k$-th entry of $\mathbf c$.
\item So a non-pivot column is the combination of pivot columns with its own entries as coefficients. Rewrite as a vector equation $\lambda_1 R_{1}+\dots+\lambda_n R_{n}=\mathbf 0$ (with a coefficient $-1$ on the non-pivot column): this is a solution of $R\boldsymbol\lambda=\mathbf 0$.
\item $A\boldsymbol\lambda=\mathbf 0$ and $R\boldsymbol\lambda=\mathbf 0$ have the same solutions, because row operations preserve solution sets (Problem 6). So the same coefficients give the same relation among the columns of $A$.
\end{enumerate}\textbf{Answer:} Each non-pivot column of $R$ equals $\sum_k c_k\,\mathbf e_k$ with $\mathbf e_k$ a pivot column; the relation is a solution of $R\boldsymbol\lambda=\mathbf 0$, hence of $A\boldsymbol\lambda=\mathbf 0$.\end{tcolorbox}
\par\vspace{2pt}
\end{document}
